\documentclass{article}
\usepackage{amsmath,amsthm,amssymb}

\newcommand{\dinv}{\operatorname{dinv}}
\newcommand{\DS}{\operatorname{DS}}
\newcommand{\DSstar}{\operatorname{DS}^{\ast}}
\newtheorem{theorem}{Theorem}
\newtheorem{conjecture}{Conjecture}
\newtheorem{definition}{Definition}
\newtheorem{remark}{Remark}

\title{Dyck Symmetric Functions}
\author{}
\date{}

\begin{document}
\maketitle

\section{Status}

This note records the status of Dyck symmetric functions in the curated
repository.  The degenerate step \(t=0\) case is classical RSK/dual-RSK
combinatorics.  The classical Dyck step \(t=1\) case is proved in the 2026
preprint
\emph{Dyck Symmetric Functions and Applications to \(q,t\)-Catalan
Polynomials}.  The proof first establishes the dual case by an explicit
tableau insertion algorithm and then derives the affine, or nondual, case
from the dual case using the standard involution on symmetric functions.

The \(r=s t+1\) analogue below is conjectural for \(t>1\).  The accompanying
code checks systematic finite parameter boxes for both the affine and dual
identities.

\section{The \(r=s t+1\) Setup}

Fix a nonnegative integer step parameter \(t\).  For any finite integer
sequence \(x=(x_0,\ldots,x_{\ell-1})\), set
\[
  \dinv_t(x)=
  \sum_{0\le i<j<\ell} d_t(x_i,x_j),
\]
where
\[
  d_t(a,b)=
  \begin{cases}
    \max(0,a+t-b),& a\le b,\\
    \max(0,b+1+t-a),& a>b.
  \end{cases}
\]
An \emph{affine rational Dyck sequence of step \(t\)} is a finite integer
sequence satisfying
\[
  x_{i+1}\le x_i+t
  \qquad\text{for all }i.
\]
A \emph{dual rational Dyck sequence of step \(t\)} is a finite integer
sequence satisfying
\[
  x_{i+1}>x_i+t
  \qquad\text{for all }i.
\]

Let \(S\) be a finite multiset of integers.  A factorization of \(S\) is a
sequence of finite words whose concatenation is a rearrangement of \(S\).  It
is affine, respectively dual, if each factor is an affine, respectively dual,
rational Dyck sequence of step \(t\).  Define
\[
  \DS^{(t)}(S,d;\mathbf x)
  =
  \sum_{\substack{\mathcal F\text{ affine factorization of }S\\
                  \dinv_t(F_0F_1F_2\cdots)=d}}
  x^{\mathcal F}
\]
and
\[
  {\DSstar}^{(t)}(S,d;\mathbf x)
  =
  \sum_{\substack{\mathcal F\text{ dual factorization of }S\\
                  \dinv_t(F_0F_1F_2\cdots)=d}}
  x^{\mathcal F}.
\]

A rational Dyck tableau of step \(t\) is a left-aligned tableau whose rows are
dual rational Dyck sequences of step \(t\), and whose columns, read from
bottom to top, are affine rational Dyck sequences of step \(t\).  Let
\(\lambda(P)\) be its shape and let \(\operatorname{RR}(P)\) be its
row-reading word.

\section{The Degenerate \(t=0\) Case}

When \(t=0\), one has \(d_0(a,b)=0\) for all entries \(a,b\), so every
rearrangement of a fixed multiset lies in the single dinv class \(d=0\).
The dual factors are strictly increasing words, while the affine factors are
weakly increasing words.  The rational Dyck tableaux of step \(0\) are exactly
row-strict semistandard tableaux in our orientation: rows are strictly
increasing and columns are weakly increasing.

Thus the step \(0\) dual identity is the classical dual RSK correspondence for
strict biwords.  Fix a multiset \(S\) and a composition
\(\alpha=(\alpha_1,\ldots,\alpha_k)\).  A dual factorization of weight
\(\alpha\) is encoded as a biword whose top row contains \(i\) repeated
\(\alpha_i\) times and whose bottom row is the corresponding factor.  The
strictness of each factor is precisely the strictness condition in the
dual-RSK input.  Dual RSK gives a bijection with pairs \((P,Q)\) of common
shape, where \(P\) is row-strict semistandard of content \(S\), and \(Q\) is
ordinary semistandard of content \(\alpha\).  Equivalently, after transposing
or reversing the usual row/column conventions, this is the standard
RSK/Kostka description of Schur coefficients.

The affine \(t=0\) identity is the corresponding ordinary RSK statement for
weakly increasing factors, with the recording tableau convention transposed by
the usual involution.  Consequently the \(t=0\) specialization is known by
classical RSK theory; it is the baseline that the \(t=1\) insertion theorem
and the \(t>1\) conjectures generalize.

\section{Classical Theorem}

\begin{theorem}[Classical dual Dyck symmetric functions]
For every finite multiset \(S\) and every \(d\ge0\), the classical dual Dyck
symmetric function satisfies
\[
  \DSstar(S,d;\mathbf x)
  =
  \sum_P s_{\lambda(P)}(\mathbf x),
\]
where \(P\) ranges over classical Dyck tableaux with entries \(S\) and
\(\dinv(\operatorname{RR}(P))=d\).
\end{theorem}

\begin{remark}
The proof in the 2026 preprint is constructive.  It gives an explicit
insertion algorithm sending dual Dyck factorizations to Dyck tableaux together
with semistandard recording data.  The local row insertion is iterated through
the rows of a tableau, and the reverse insertion proves bijectivity.
\end{remark}

\begin{theorem}[Classical affine Dyck symmetric functions]
For every finite multiset \(S\) and every \(d\ge0\), the classical affine Dyck
symmetric function satisfies
\[
  \DS(S,d;\mathbf x)
  =
  \sum_P s_{\lambda(P)'}(\mathbf x),
\]
where the indexing set \(P\) is the same set of classical Dyck tableaux.
\end{theorem}

\begin{remark}
In the preprint, this nondual statement is derived from the dual statement by
comparing the fundamental-quasisymmetric expansions and applying the standard
involution \(\omega\), which sends \(s_\lambda\) to \(s_{\lambda'}\).
\end{remark}

\section{Conjectural \(r=s t+1\) Analogue}

\begin{conjecture}[Rational Dyck symmetric functions, \(r=s t+1\)]
Let \(S\) be a finite multiset of integers, let \(d\ge0\), and fix \(t\ge0\).
Then
\[
  \DS^{(t)}(S,d;\mathbf x)
  =
  \sum_P s_{\lambda(P)'}(\mathbf x),
\]
where \(P\) ranges over rational Dyck tableaux of step \(t\) whose entries
are exactly \(S\) and whose row-reading word satisfies
\[
  \dinv_t(\operatorname{RR}(P))=d.
\]
Similarly,
\[
  {\DSstar}^{(t)}(S,d;\mathbf x)
  =
  \sum_P s_{\lambda(P)}(\mathbf x),
\]
with the same indexing set of rational Dyck tableaux.
\end{conjecture}

\begin{remark}
The displayed statement is known for \(t=0\) by the classical RSK and dual
RSK correspondences, and for \(t=1\) by the 2026 Dyck insertion theorem.  For
\(t>1\), it is currently treated as a conjecture.  The directory
\texttt{code/} includes systematic finite checks over bounded multisets and
all occurring dinv values for the following official parameter boxes:
\[
  (t,A,L)=(2,10,10),\qquad (3,13,9),\qquad (4,16,8),
\]
where \(A\) is the alphabet size \(\{1,\ldots,A\}\) and \(L\) is the maximum
word length.  The same directory also includes Monte Carlo class checks: each
trial samples a word uniformly from \(\{1,\ldots,A\}^L\), fixes its multiset
and dinv value, and then exhaustively checks that sampled class.  The official
Monte Carlo runs use \(100\) sampled classes for
\[
  (t,A,L)=(2,11,12),\qquad (3,14,11),\qquad (4,17,10).
\]
The directory also includes the classical insertion algorithm code used for
the theorem-level case.
\end{remark}

\end{document}
